\documentclass[conference]{IEEEtran}
\usepackage{graphicx}
\usepackage{amsmath}
\usepackage{amssymb}
\usepackage{booktabs}
\usepackage{cite}
\usepackage{float}
\usepackage{subcaption}
\usepackage{array}
\usepackage{multirow}
\usepackage{xcolor}
\usepackage{hyperref}

\title{Time Series Forecasting Using Recurrent Neural Networks:
A Comparative Analysis of Elman, Jordan, and Multi-Recurrent Architectures}

\author{
    \IEEEauthorblockN{M Tashreeq Hartogh}
    \IEEEauthorblockA{
        27163245\\
        Department of Computer Science\\
        Stellenbosch University\\
        Stellenbosch, South Africa\\
        https://github.com/calebwidogast96-art/RNN-Forecasting
    }
}

\begin{document}

\maketitle

\begin{abstract}
This report presents a comprehensive comparative analysis of three
recurrent neural network (RNN) architectures for time series forecasting:
Elman RNNs (Simple RNN), Jordan RNNs, and Multi-Recurrent RNNs. I evaluate
these models on five benchmark datasets (Electricity, ETTh1, COVID-19,
Weather, and Traffic) that exhibit varying degrees of stationarity
and complexity. Using a thorough empirical methodology incorporating
time series cross-validation, data preprocessing, and systematic
hyperparameter tuning, I assess model performance across different lookback
windows and hidden unit configurations. The results show that
Multi-Recurrent RNNs with multi-scale temporal dependencies achieved
superior performance on most datasets, with Mean Absolute Error (MAE)
values consistently lower than traditional Elman RNNs. This study provides
insights into the effectiveness of different RNN architectures for practical
time series forecasting applications and discusses the trade-offs between model
complexity and predictive accuracy.
\end{abstract}

\begin{IEEEkeywords}
time series forecasting, recurrent neural networks, Elman RNN,
Jordan RNN, multi-recurrent RNN, deep learning
\end{IEEEkeywords}

% ============================================================
% INTRODUCTION
% ============================================================

\section{Introduction}

Time series forecasting is a fundamental problem in machine learning
with applications in finance, weather prediction, energy demand
forecasting, epidemiology, and traffic management. Unlike traditional
supervised learning where samples are independent and identically distributed
(i.i.d.), time series data contains temporal dependencies and correlated
structures that require specialized modeling approaches.

Recurrent Neural Networks (RNNs) have emerged as a powerful class of
models for time series tasks due to their ability to maintain hidden state
across time steps, effectively capturing temporal patterns and dependencies.
However, the variety of RNN architectures available presents practitioners
with important design choices that can significantly affect forecasting
accuracy.

\subsection{Motivation and Research Questions}

The main aim for this study is to systematically compare
and evaluate different RNN architectures designed to capture temporal
information at varying scales. Specifically, I investigate the following
research questions:

\begin{enumerate}
    \item How do different RNN architectures (Elman, Jordan,
    Multi-Recurrent) perform across diverse time series datasets with
    different statistical properties?
    
    \item What is the impact of lookback window size and hidden unit count
    on model performance?
    
    \item Which architectural features (e.g., output feedback, multi-scale
    delays) provide the most benefit for time series forecasting?
    
    \item How do these architectural innovations compare to established
    baselines (GRU, LSTM)?
\end{enumerate}

\subsection{Contributions}

This work makes the following contributions:

\begin{itemize}
    \item An empirical evaluation of RNN variants using
    thorough time series validation methodology.
    
    \item Implementation of a Multi-Recurrent RNN architecture
    with configurable multi-scale temporal dependencies.
    
    \item Comprehensive analysis of hyperparameter effects (lookback,
    hidden units) across multiple datasets.
    
    \item Clear guidelines for practitioners on architecture selection
    for time series forecasting.
\end{itemize}

\subsection{Report Structure}

The remainder of this report is organized as follows:

Section~\ref{sec:background} provides theoretical background on RNN
architectures and time series methodologies.

Section~\ref{sec:datasets} describes the five datasets used in my 
analysis.

Section~\ref{sec:implementation} details data preprocessing and
implementation specifics.

Section~\ref{sec:empirical} describes my empirical process and
experimental methodology.

Section~\ref{sec:results} presents and discusses the results. Finally,

Section~\ref{sec:conclusions} concludes with recommendations for
practitioners.

% ============================================================
% BACKGROUND
% ============================================================

\section{Background}
\label{sec:background}

\subsection{Recurrent Neural Networks}

Recurrent Neural Networks are a class of neural networks designed to
process sequential data by maintaining a hidden state $h_t$ that is
updated at each time step. The fundamental recurrence relation is:

\begin{equation}
h_t = \sigma(W_{hh}h_{t-1} + W_{xh}x_t + b_h)
\end{equation}

where $x_t$ is the input at time step $t$, $W_{hh}$ and $W_{xh}$ are
weight matrices, $b_h$ is the bias term, and $\sigma$ is an
activation function (typically tanh or ReLU).

\subsubsection{Elman RNN (Simple RNN)}

The Elman RNN, also known as a Simple RNN, is the basic RNN architecture
where the hidden state from the previous time step is fed back as input
to the current step. The output at time $t$ is:

\begin{equation}
y_t = W_{hy}h_t + b_y
\end{equation}

While seemingly simple, Elman RNNs suffer from vanishing and exploding
gradient problems during backpropagation through time (BPTT), limiting
their ability to capture long-term dependencies.

\subsubsection{Jordan RNN}

The Jordan RNN improves upon the Elman architecture by replacing the
hidden state feedback with output feedback. The recurrence is:

\begin{equation}
h_t = \sigma(W_{hx}x_t + W_{yh}y_{t-1} + b_h)
\end{equation}

where $y_{t-1}$ is the output from the previous time step, not the hidden
state. This design has the advantage of a clearer interpretation of
what is being fed back (actual predictions) and can reduce the gradient
problem by providing a shorter error propagation path.

\subsubsection{Multi-Recurrent RNN}

The Multi-Recurrent RNN extends the standard architecture by using
connections from multiple previous hidden states at different delay intervals:

\begin{equation}
h_t =
\sigma\left(
W_{hx}x_t +
\sum_{d \in \text{delays}} W_d h_{t-d} +
b_h
\right)
\end{equation}

This architecture explicitly models temporal dependencies at multiple
scales, allowing the network to learn patterns at different frequencies.
For instance, with delays $(1,6,12)$, the network can capture immediate effects,
weekly patterns, and monthly patterns simultaneously.
\\
\subsection{Preventing Overfitting and Underfitting}

\subsubsection{Regularization Strategies}

To prevent overfitting while avoiding underfitting, I implement:

\begin{itemize}
    \item \textbf{Early Stopping}: Monitor validation loss and stop
    training when it increases for 5 consecutive epochs, restoring
    the best weights. This prevents overfitting while ensuring the
    model has trained sufficiently.
    
    \item \textbf{Validation Set Monitoring}: Use separate validation
    and test sets to distinguish between underfitting and overfitting.
\end{itemize}

\subsubsection{Capacity Control}

Model capacity is controlled through:

\begin{itemize}
    \item Systematic variation of hidden units (8 to 512) to find the
    optimal complexity.
    
    \item Careful selection of lookback window to avoid both too-short
    contexts (underfitting) and too-long contexts (overfitting and
    computational inefficiency).
\end{itemize}

\subsection{Performance Metrics}

\subsubsection{Mean Squared Error (MSE)}

The loss function used for training is:

\begin{equation}
\text{MSE} =
\frac{1}{N}\sum_{i=1}^{N}(y_i - \hat{y}_i)^2
\end{equation}

MSE penalizes larger errors more heavily, making it sensitive to
outliers.

\subsubsection{Mean Absolute Error (MAE)}

Reported as the primary evaluation metric:

\begin{equation}
\text{MAE} =
\frac{1}{N}\sum_{i=1}^{N}|y_i - \hat{y}_i|
\end{equation}

MAE is scale-dependent but interpretable in the original data units
and more robust to outliers than MSE. To account for the scale of the
data, min-max scaling was applied. Thus all MAE values throughout
the rest of this report are scaled.

% ============================================================
% DATASETS
% ============================================================
\section{Datasets}
\label{sec:datasets}

We selected five diverse benchmark datasets that exhibit different temporal characteristics, allowing comprehensive evaluation of RNN architectures.

\subsection{Dataset Overview}

\begin{table}[H]
\centering
\caption{Summary of datasets}
\label{tab:datasets}
\begin{tabular}{lcccc}
\toprule
\textbf{Dataset} & \textbf{Variables} & \textbf{Granularity} &
\textbf{Length} & \textbf{Stationarity} \\
\midrule
Electricity & 321 & 15 minutes & 26,304 obs & Non-stationary \\
ETTh1 & 6 & 1 hour & 17,420 obs & Non-stationary \\
Weather & 21 & 10 minutes & 52,560 obs & Non-stationary \\
Traffic & 430 & 1 hour & 17,544 obs & Non-stationary \\
COVID-19 & 55 & 1 day & 365 obs & Non-stationary \\
\bottomrule
\end{tabular}
\end{table}

\subsection{Electricity Load Dataset}

\textbf{Source:} UCI Machine Learning Repository
(ElectricityLoadDiagrams20112014).

The electricity dataset contains electricity consumption data from
321 individual clients measured at 15-minute intervals over several
years. The target variable is the origin (OT) electricity consumption.

\textbf{Characteristics:}

\begin{itemize}
    \item High-dimensional multivariate time series (320 features +
    1 target).
    \item Clear daily and weekly seasonality patterns.
    \item Non-stationary due to trend and seasonal components.
\end{itemize}

\subsection{ETT Dataset (Electricity Transformer Temperature)}

\textbf{Source:} GitHub ETDataset Repository.

The Electricity Transformer Temperature (ETT) dataset comprises
measurements from electric power transformers
including six load variables: HUFL (High Use Ful Load),
HULL (High Use Low Load), MUFL (Middle Use Ful Load),
MULL (Middle Use Low Load), LUFL (Low Use Ful Load), and
LULL (Low Use Low Load).

\textbf{Characteristics:}

\begin{itemize}
    \item Moderate dimensionality (6 variables).
    \item Hourly granularity captures daily heating cycles.
    \item Clear periodicity and trend components.
    \item Relatively clean data with minimal missing values.
\end{itemize}

\subsection{Weather Dataset}

\textbf{Source:} Max Planck Institute for Biogeochemistry.

This dataset contains 21 meteorological indicators collected at
10-minute intervals, including temperature, humidity, pressure,
radiation, and precipitation measurements.

\textbf{Characteristics:}

\begin{itemize}
    \item Moderate dimensionality (21 variables).
    \item High-frequency sampling (10-minute intervals).
    \item Strong seasonal patterns (daily and annual cycles).
    \item Non-stationary trend due to seasonal changes.
\end{itemize}

\subsection{Traffic Dataset}

\textbf{Source:} Traffic Data Repository.

The traffic dataset comprises hourly traffic speed measurements from
430 San Francisco freeway car lanes.

\textbf{Characteristics:}

\begin{itemize}
    \item High-dimensional (430 variables representing different free lanes).
    \item Clear weekly patterns (weekday vs. weekend).
    \item Multiple peaks during rush hours.
    \item Non-stationary due to urban development and seasonal variations.
\end{itemize}

\subsection{COVID-19 Dataset}

\textbf{Source:} Johns Hopkins University CSSE.

Contains COVID-19 hospitalization data for the U.S. state of
California (CA), provided by Johns Hopkins University.

\textbf{Characteristics:}

\begin{itemize}
    \item Moderate dimensionality (55 series).
    \item Daily sampling (lowest frequency in our collection).
    \item Non-stationary with multiple structural breaks (lockdowns, policy changes, variants).
    \item Highly irregular patterns not amenable to simple seasonality.
\end{itemize}

% ============================================================
% IMPLEMENTATION
% ============================================================

\section{Implementation}
\label{sec:implementation}

\subsection{Data Preprocessing}

All datasets undergo identical preprocessing to ensure fair
comparison.

\subsubsection{Train/Validation/Test Split}

Chronological splitting is applied to respect temporal ordering:

\begin{itemize}
    \item \textbf{Training Set}: First 70\% of chronological data.
    \item \textbf{Validation Set}: Next 15\% (between 70\% and 85\%).
    \item \textbf{Test Set}: Final 15\% (after 85\%).
\end{itemize}

This split simulates realistic scenarios where models are trained on
historical data and evaluated on unseen future data.

\subsubsection{Normalization}

Min-Max scaling is applied to each variable independently:

\begin{equation}
x_{\text{scaled}} =
\frac{x - \min(x)}
{\max(x) - \min(x)}
\end{equation}

where min and max are computed on the training set only, then applied
to validation and test sets to prevent information leakage.

\subsubsection{Sequence Creation}

Time series are converted into supervised learning problems using a
sliding window approach. For a given lookback window $L$, we create
sequences:

\begin{equation}
\left\{
(x_{t-L}, \ldots, x_{t-1}, y_t)
:
t = L+1, \ldots, T
\right\}
\end{equation}

where $x_t$ represents the multivariate features at time $t$ and
$y_t$ is the target variable at time $t$.

\subsection{Model Architecture}

\subsubsection{Network Configuration}

All RNN variants are configured as follows:

\begin{itemize}
    \item \textbf{Input Layer}: Accepts sequences of shape
    $(\text{batch\_size}, \text{lookback}, \text{num\_features})$.
    
    \item \textbf{RNN Layer}: Variable hidden units from
    $\{8,16,32,64,128,256,512\}$.
    
    \item \textbf{Output Layer}: Single neuron with linear activation for regression.
    
    \item \textbf{Activation Function}: Hyperbolic tangent (tanh) for hidden units.
\end{itemize}

\subsubsection{Optimization}

\begin{itemize}
    \item \textbf{Optimizer}: Adam (Adaptive Moment Estimation).
    \item \textbf{Learning Rate}: $10^{-3}$ (default).
    \item \textbf{Loss Function}: Mean Squared Error (MSE).
    \item \textbf{Batch Size}: 32 samples.
    \item \textbf{Maximum Epochs}: 50.
    \item \textbf{Early Stopping}: Patience of 5 epochs on
    validation loss.
\end{itemize}

\subsection{Experimental Framework}

Implementation is conducted in Python 3.x using TensorFlow 2.x with
custom layer implementations for Jordan and Multi-Recurrent RNNs to
enable precise control over recurrent connections.
(See \textbf{custom\_layers.py} in the GitHub repository)

\begin{itemize}
    \item \textbf{JordanRNN}: Custom layer implementing output
    feedback mechanism.
    \item \textbf{MultiRecurrentRNN}: Custom layer implementing
    multi-scale delays.
    \item \textbf{SimpleRNN}: TensorFlow's built-in Elman RNN layer.
\end{itemize}

% ============================================================
% EMPIRICAL PROCESS
% ============================================================

\section{Empirical Process}
\label{sec:empirical}

\subsection{Experimental Design}

The empirical process follows a systematic grid search over
hyperparameters, each architecture with differing amount of 
hidden units at differing loockback ranges while maintaining
temporal integrity.

\subsubsection{Hyperparameter Grid}

For each dataset and architecture, we systematically vary:

\begin{table}[H]
\centering
\caption{Hyperparameter search grid}
\label{tab:hyperparams}
\begin{tabular}{lccc}
\toprule
\textbf{Architecture} & \textbf{Hidden Units} \\
\midrule
Simple RNN & 8, 16, 32, 64, 128, 256, 512 \\
Jordan RNN & 8, 16, 32, 64, 128, 256 \\
Multi-Recurrent RNN & 64 (fixed) \\
\bottomrule
\end{tabular}
\end{table}

Where Multi-Recurrent RNN was tested across the following delays:
$(1)$, $(1,2)$, $(1,2,3)$, $(1,2,3,4)$,
$(1,6,12)$, $(1,3,6,12)$, $(1,6,12,24)$ 
\\
\textbf{Dataset-Specific Lookback Ranges:}

\begin{itemize}
    \item \textbf{Electricity}: 6, 12, 18, 24, 30, 36, 42, 48 (6-hour to 48-hour intervals).
    
    \item \textbf{ETTh1}: 6, 12, 18, 24, 30, 36, 42, 48, 54, 60 (6-hour to 60-hour).
    
    \item \textbf{Weather}: 12, 24, 36 (2-hour to 6-hour windows at 10-min granularity).
    
    \item \textbf{Traffic}: 12, 24, 36, 48 (12-hour to 48-hour windows).
    
    \item \textbf{COVID-19}: 1 through 21 (1 to 21 days).
\end{itemize}

\subsection{Experimental Procedure}

\begin{enumerate}
    \item \textbf{Data Preparation}: Load raw dataset and apply
    train/validation/test split respecting temporal order.
    
    \item \textbf{Preprocessing}: Scale using MinMax normalization
    (fit on training set), then create sequences with the specified
    lookback.
    
    \item \textbf{Model Training}: For each combination of
    (architecture, lookback, hidden units):
    \begin{enumerate}
        \item Initialize model with random weights.
        \item Train on training set using Adam optimizer and MSE
        loss.
        \item Monitor validation loss and apply early stopping
        (patience = 5).
        \item Restore best weights from validation phase.
    \end{enumerate}
    
    \item \textbf{Evaluation}: Compute MAE and MSE on held-out test
    set using denormalized predictions.
    
    \item \textbf{Results Storage}: Record all hyperparameters and
    metrics to JSON for subsequent analysis.
\end{enumerate}

\subsection{Justification of Methodological Choices}

\subsubsection{Why Chronological splitting Validation?}

Chronological splitting validation respects the temporal structure
of time series data, preventing information leakage from future to
past. Random shuffling or standard k-fold cross-validation would
violate the fundamental assumption that future values can not
influence past values.

\subsubsection{Why MinMax Scaling?}

MinMax scaling bounds values to $[0,1]$, which:

\begin{itemize}
    \item Improves neural network training stability.
    \item Allows fair comparison across datasets with
    different scales.
    \item Preserves the distribution shape (unlike standardization).
\end{itemize}

\subsubsection{Why Early Stopping with Patience=5?}

This configuration prevents overfitting by stopping when validation
loss plateaus while allowing sufficient training. A patience of 5
epochs provides stability without premature termination.

\subsection{Computational Considerations}

Total experimental runs: Approximately 1,500+ model trainings
(varying by dataset). Each model trains for typically 20--40 epochs
before early stopping. All experiments use GPU acceleration where
available.

% ============================================================
% RESULTS
% ============================================================

\section{Results and Discussion}
\label{sec:results}

\subsection{Overall Performance Summary}

We present comprehensive results across all five datasets, comparing three RNN architectures (Simple RNN, Jordan RNN, Multi-Recurrent RNN) against baseline methods (GRU, LSTM).

\subsubsection{Best Performance by Dataset}

\begin{table}[H]
\centering
\caption{Best test MAE achieved for each dataset and architecture}
\label{tab:best_results}
\begin{tabular}{lcccc}
\toprule
\textbf{Dataset} & \textbf{Architecture} & \textbf{Test MAE} &
\textbf{H. Units/Delay} & \textbf{Lookback} \\
\midrule
Electricity & GRU & 0.0286 & 128 & 30 \\
ETTh1 & Multi-Recurrent & 0.0579 & (1, 2, 3) & 60 \\
Weather & Multi-Recurrent & 0.0011 & (1, 6, 12) & 36 \\
Traffic & GRU & 0.0279 & 64 & 36 \\
COVID-19 & Jordan & 0.1848 & 128 & 3 \\
\bottomrule
\end{tabular}
\end{table}

Multi-Recurrent RNN performed best on ETTh1 and Weather, which
contain 6 and 21 input features, respectively. GRU performed best on
Electricity and Traffic, which contain 321 and 430 input features,
respectively. In contrast, Jordan performed best on the sparse
COVID-19 dataset, which contains only 365 observations. Overall, the
results suggest that model performance varies across
datasets, with the best recurrent architecture depending on
the characteristics and dimensionality of the data.

\subsection{Electricity Dataset Results}

\subsubsection{Architecture Comparison}

Among the evaluated architectures, the GRU achieved the best
performance with 128 units and a lookback of 30 observations,
corresponding to approximately 7.5 hours of historical data.
GRUs provide gated memory that can retain relevant information over
time while using fewer parameters than LSTMs, making them well suited
to capturing recurring consumption patterns without unnecessary model
complexity. This likely contributed to their advantage over the simpler
RNN, while the additional complexity of LSTMs and multi-recurrent
architectures did not provide a corresponding improvement. The Jordan
architecture was also less suited to this task than the GRU's direct
gated handling of temporal dependencies.


\subsubsection{Effect of Lookback Window}

\begin{table}[H]
\centering
\caption{Electricity dataset: Test MAE for Multi-Recurrent RNN across lookback windows}
\label{tab:electricity_lookback}
\begin{tabular}{cc|cc}
\toprule
Lookback & MAE & Lookback & MAE \\
\midrule
6 & 0.0316 & 30 & 0.0331 \\
12 & 0.0300 & 36 & 0.0315 \\
18 & 0.0329 & 42 & 0.0315 \\
24 & 0.0317 & 48 & 0.0318 \\
\bottomrule
\end{tabular}
\end{table}

Lookback values between 12--36 perform best, capturing both
immediate correlations and daily patterns (electricity consumption
has 96 measurements per day at 15-min granularity). Very short
windows (6) underfit, while very long windows (42--48) suffer from
overfitting.

\subsection{ETTh1 Dataset Results}

The ETT dataset shows clearer patterns due to the transformer
heating cycle and lower dimensionality (6 features).

\subsubsection{Architecture Comparison}

\begin{table}[H]
\centering
\caption{ETTh1: Test MAE for each architecture using optimal
hyperparameters}
\label{tab:etth1_architecture}
\begin{tabular}{cc|cc}
\toprule
\textbf{Architecture} & \textbf{Best MAE} & \textbf{Architecture} & \textbf{Best MAE} \\
\midrule
Simple RNN & 0.1199 & GRU & 0.1082 \\
Jordan RNN & 0.1535 & LSTM & 0.1123 \\
Multi-Rec & 0.0857 & Improvement & 20.7\% vs GRU \\
\bottomrule
\end{tabular}
\end{table}

The Multi-Recurrent RNN achieved the best performance on the ETTh1 dataset,
with a test MAE of 0.0857, representing a 28.4\% improvement over the Simple RNN
and a 20.7\% improvement over the GRU. This suggests that the multiple recurrent
pathways were better able to capture the different temporal scales present in the
transformer temperature data. In particular, the delayed recurrent connections
allow the model to incorporate both short-term changes in temperature and longer-term
patterns associated with hourly and daily heating cycles. The GRU and LSTM also
benefited from their gated memory mechanisms, achieving MAEs of 0.1082 and 0.1123
respectively, but did not match the Multi-Recurrent architecture. The Jordan RNN
performed worst, with an MAE of 0.1535, suggesting that feeding previous outputs back
into the network was less effective for capturing the temporal dependencies in this
dataset.

\subsection{Weather Dataset Results}

The weather dataset, with highest sampling frequency (10 minutes) 
and clean patterns, showed the smallest absolute errors.

\subsubsection{Performance Table by Architecture and Lookback}

\begin{table}[H]
\centering
\caption{Weather dataset: Test MAE across architectures and key lookback values}
\label{tab:weather_results}
\begin{tabular}{llll}
\toprule
\textbf{Architecture} & Lookback 12 & Lookback 24 & Lookback 36 \\
\midrule
Simple RNN (128) & 0.0028 & 0.0016 & 0.0015 \\
Jordan RNN (128) & 0.0030 & 0.0018 & 0.0016 \\
Multi-Recurrent (64) & 0.0017 & 0.0012 & 0.0012 \\
GRU (128) & 0.0022 & 0.0016 & 0.0015 \\
LSTM (128) & 0.0020 & 0.0014 & 0.0013 \\
\bottomrule
\end{tabular}
\end{table}

The weather dataset exhibits strong periodic behaviour, which makes
it well suited to architectures capable of modelling temporal dependencies
at multiple time scales. The Multi-Recurrent RNN consistently achieved the lowest
MAE across the tested lookback windows, reaching 0.0012 at both 24 and 36 lookback.
Its delayed recurrent connections allow the model to incorporate information from
different points in the recent history, which is useful for capturing both short-term
fluctuations and recurring environmental patterns. The LSTM also performed well,
achieving an MAE of 0.0013 at a lookback of 36, while the GRU and Simple RNN achieved
0.0015. The Jordan RNN produced the highest errors across all three lookback values,
suggesting that its output-feedback mechanism was less effective for this type of
highly periodic data. Overall, the consistently low errors across the larger lookback
windows indicate that the weather series benefits from retaining a broader temporal
context, with the Multi-Recurrent architecture providing the strongest performance among
the evaluated models.

\subsection{Traffic Dataset Results}

Traffic forecasting presents particular challenges due to irregular patterns and external events.

\subsubsection{Architecture Comparison}

The reported test MAE values for the Traffic dataset are:

\begin{table}[H]
\centering
\caption{Traffic dataset: Test MAE by architecture}
\label{tab:traffic_results}
\begin{tabular}{lc}
\toprule
\textbf{Architecture} & \textbf{Test MAE} \\
\midrule
GRU & 0.0279 \\
Multi-Recurrent & 0.0327 \\
LSTM & 0.0343 \\
Simple RNN & 0.0387 \\
Jordan RNN & 0.0369 \\
\bottomrule
\end{tabular}
\end{table}

The GRU achieved the lowest test MAE of 0.0279 using a lookback of 36
and 64 units. This was followed by the Multi-Recurrent RNN with an MAE
of 0.0327, while the LSTM, Jordan RNN, and Simple RNN achieved 0.0343,
0.0369, and 0.0387 respectively. The stronger performance of the GRU
suggests that its gated memory mechanism was better suited to the irregular
temporal dependencies present in traffic data. Unlike highly periodic datasets
such as weather or electricity consumption, traffic patterns can vary substantially
due to congestion, incidents, weather, and other external factors. The
GRU gates allow the model to selectively retain or discard information as
these patterns change, providing greater flexibility than the fixed delayed
connections of the Multi-Recurrent architecture. The LSTM also provides gated
memory, but its additional complexity did not translate into improved performance
in this case. The poorer performance of the Simple and Jordan RNNs further
suggests that simpler recurrent mechanisms were less effective at handling the
variable temporal dependencies in the traffic series.

\subsection{COVID-19 Dataset Results}

COVID-19 data presented the most challenging forecasting scenario due to non-stationarity and structural breaks.

\subsubsection{Performance Patterns}

\begin{table}[H]
\centering
\caption{COVID-19: Best Test MAE for Jordan RNN per lookback values}
\label{tab:covid_lookback}
\begin{tabular}{cc|cc}
\toprule
Lookback & MAE & Lookback & MAE \\
\midrule
1 & 0.3915 & 6 & 0.4081 \\
2 & 0.2621 & 7 & 0.4347 \\
3 & 0.1848 & 8 & 0.5838 \\
4 & 0.4586 & 9 & 0.5163 \\
5 & 0.7425 & 10 & 0.8279 \\
\bottomrule
\end{tabular}
\end{table}

The Jordan RNN achieved its lowest MAE of 0.1848 at a lookback of 3, after which performance
deteriorated considerably. In particular, increasing the lookback beyond 3
did not provide additional useful historical information, with the MAE increasing
to 0.7425 at a lookback of 5 and reaching 0.8279 at a lookback of 10, with similar
degrading values up to the lookback of 21. This behaviour is consistent with the highly non-stationary nature of COVID-19
case data, where changes in transmission, interventions, reporting practices,
and other external factors can cause historical patterns to become less representative
of future observations.
\\
Compared with the other architectures, the Jordan RNN was particularly sensitive
to the amount of historical information provided. Its reliance on previous outputs
may have allowed it to respond effectively to the very short-term dependencies in
the series at a lookback of 3, but the additional historical context appears to have
introduced information that was less useful for forecasting. In contrast, architectures
such as the GRU and LSTM use gated memory mechanisms to selectively retain relevant
information, which can provide greater control over how historical observations influence
the prediction. The Multi-Recurrent architecture similarly provides access to information
at multiple temporal delays, but the usefulness of these delays is limited when the underlying
process contains frequent structural changes. Overall, the results suggest that short-term
information was more valuable than long historical windows for this dataset, with the sharp
variation in performance across lookback values highlighting the difficulty of applying recurrent
models to non-stationary time series.

\subsection{Cross-Dataset Analysis}

\subsubsection{Architecture Effectiveness}

Overall, the results show that no single recurrent architecture consistently
performed best across all datasets. Performance depended strongly on the
temporal characteristics of each series. The Multi-Recurrent RNN 
particularly well on the more periodic ETTh1 and Weather datasets,
while the GRU achieved the lowest error on the more irregular Traffic dataset.
The Electricity dataset also favoured recurrent models with sufficient temporal
memory, with the GRU performing best among the evaluated architectures.
The COVID-19 dataset was substantially more difficult, with strong sensitivity
to lookback length due to its non-stationary behaviour and structural changes.
These results indicate that model architecture and lookback length should be selected
according to the underlying temporal structure of the forecasting problem rather
than assuming that one architecture will generalise best across all time series.

\subsubsection{Hidden Unit Counts}

\begin{table}[H]
\centering
\caption{Observed performance by hidden unit count}
\label{tab:hidden_units_trend}
\begin{tabular}{lc}
\toprule
\textbf{Hidden Units} & \textbf{Observed Performance} \\
\midrule
8 & Good \\
16 & Good \\
32 & Better \\
64 & Best \\
128 & Good \\
256 & Worse \\
512 & Poor \\
\bottomrule
\end{tabular}
\end{table}

The results suggest that increasing the number of hidden units does
not necessarily improve forecasting performance. Moderate model sizes,
particularly 64 and 128 units, generally provided a good balance between
model capacity and generalisation, while substantially larger configurations
did not consistently produce improvements. Smaller configurations of 8--16
units may have limited capacity to represent more complex temporal dependencies,
whereas increasing the number of units beyond the moderate range introduces
additional parameters without a corresponding improvement in test performance.
The results therefore suggest that moderate hidden-layer sizes are generally
sufficient for these datasets, although the optimal number of units remains
dependent on the architecture and characteristics of the individual time series.


\subsection{Discussion of Key Findings}

\subsubsection{Architecture Performance}

The results demonstrate that no single recurrent architecture
consistently outperformed the others across all datasets. Instead,
performance depended strongly on the temporal characteristics of the
forecasting problem. Multi-Recurrent RNNs performed particularly well
on the ETTh1 and Weather datasets, which exhibited clearer periodic
behaviour and temporal patterns at multiple scales. In contrast, the
GRU achieved the lowest reported MAE on the Traffic dataset,
suggesting that its gated memory mechanism was better suited to the more
irregular temporal dependencies present in that series. The Electricity
results also showed that relatively complex recurrent architectures
could effectively model consumption patterns, while the COVID-19
results demonstrated substantially greater sensitivity to the choice
of lookback window.

\subsubsection{Dataset-Dependent Observations}

The results highlight clear differences between datasets. The
Weather and ETTh1 datasets exhibited strong periodic patterns and
generally benefited from architectures capable of representing temporal
dependencies across multiple scales. The Traffic dataset was more
irregular, for which the GRU achieved the lowest reported MAE of 0.0279.
The Electricity dataset contained substantial temporal structure due to
its high-frequency sampling and recurring consumption behaviour, with
moderate lookback windows generally performing well. In contrast, the
COVID-19 dataset was considerably more difficult to forecast, with
large changes in performance across lookback values reflecting its non-stationary
behaviour and structural changes.

\subsubsection{Effect of Lookback}

Lookback length had a substantial effect on model performance,
although the optimal value varied between datasets. For the
Electricity dataset, lookbacks between 12 and 36 observations performed
relatively consistently, with the best reported MAE of 0.0300 occurring
at 12 observations. For the Weather dataset, performance improved as the
lookback increased from 12 to 24--36 observations. In contrast, the
COVID-19 dataset favoured a very short lookback, with the Jordan RNN
achieving its lowest reported MAE at only 3 observations. These results
indicate that the amount of historical information required is
strongly dependent on the temporal structure and stationarity of the
underlying series.

\subsubsection{Practical Implications}

The experiments suggest that recurrent architecture and
hyperparameter selection should be treated as dataset-specific
decisions rather than assuming a universally optimal model.
Architectures with multi-scale temporal connections can be
particularly effective for strongly periodic series, while gated
architectures such as GRUs may be more suitable when temporal
dependencies are irregular. Similarly, hidden-unit count and
lookback length should be tuned together with the architecture, as
increasing model capacity or historical context does not necessarily
improve generalisation. Overall, the results emphasise the importance
of matching the model's temporal structure to the characteristics
of the forecasting dataset.


\section{Conclusions}

\label{sec:conclusions}

\subsection{Summary of Findings}

This empirical study evaluated five recurrent architectures across five time series datasets with diverse temporal characteristics. The results demonstrate:

\begin{enumerate}

    \item \textbf{Architecture performance is dataset-dependent}: no
    single architecture consistently achieved the lowest test
    MAE across all datasets.

    \item \textbf{Multi-Recurrent RNNs performed strongly on periodic
    data}, achieving the lowest reported MAE on the ETTh1 and
    Weather datasets, while GRU performed best on the Traffic dataset.

    \item \textbf{Model capacity should be moderate}: configurations
    using 64--128 hidden units generally performed well, while increasing
    the number of units did not consistently improve generalization.

    \item \textbf{Lookback selection is important}: the optimal historical
    context varied substantially between datasets, with shorter windows
    favoured for highly non-stationary data and longer windows benefiting
    datasets with clearer temporal patterns.

\end{enumerate}
\newpage
\subsection{Contributions}

This work contributes:

\begin{itemize}
    \item A systematic comparison of Multi-Recurrent RNNs with conventional recurrent architectures on benchmark datasets.
    
    \item An empirical investigation of multi-scale temporal modeling in RNNs.
    
    \item Analysis of architecture and hyperparameter effects across multiple datasets.
    
    \item A time series evaluation methodology designed to avoid temporal information leakage.
\end{itemize}

\subsection{Limitations and Future Work}

\subsubsection{Limitations}

\begin{itemize}
    \item The study is limited to single-step-ahead forecasting; multi-step horizons warrant further investigation.
    
    \item Fixed delay patterns are used in the Multi-Recurrent RNN; learnable delays could provide additional flexibility.
    
    \item Computational costs were not thoroughly analyzed; GRU and LSTM may have different computational characteristics depending on implementation.
\end{itemize}

\subsubsection{Future Directions}

\begin{enumerate}
    \item Investigate adaptive delays in Multi-Recurrent RNNs.
    \item Combine gating mechanisms (GRU) with multi-scale modeling.
    \item Apply attention mechanisms to weight delay importance.
    \item Extend the evaluation to multi-step forecasting and probabilistic predictions.
    \item Evaluate the architectures on additional domains such as finance and speech.
\end{enumerate}

% ============================================================
% REFERENCES
% ============================================================

\begin{thebibliography}{99}

\bibitem{RNNForecasting}

M. Tashreeq Hartogh,

\textit{RNN Forecasting},

GitHub repository.

Available: \url{https://github.com/calebwidogast96-art/RNN-Forecasting}

\bibitem{JordanGithub}

M. Ebadpour,

\textit{Recurrent Neural Networks (RNNs) and Implementing Elman and Jordan Architectures},

GitHub repository, 2024.

Available: \url{https://github.com/MohsenEbadpour/Recurrent-Neural-Networks--RNNs--and-implementing-Elman-and-Jordan-architectures}


\bibitem{KerasRNN}

Keras,

``RNN layer,''

\textit{Keras API Documentation}.

Available: \url{https://keras.io/api/layers/recurrent_layers/rnn/}


\bibitem{MultivariateTimeSeries}

L. Lai,

\textit{Multivariate Time Series Data},

GitHub repository.

Available: \url{https://github.com/laiguokun/multivariate-time-series-data}


\bibitem{ElectricityUCI}

A. Trindade,

``ElectricityLoadDiagrams20112014,''

\textit{UCI Machine Learning Repository}, 2015.

Available: \url{https://archive.ics.uci.edu/dataset/321/electricityloaddiagrams20112014}


\bibitem{WeatherMPI}

Max Planck Institute for Biogeochemistry,

``Weather Station,''

Jena, Germany.

Available: \url{https://www.bgc-jena.mpg.de/wetter/}


\bibitem{ETDataset}

H. Zhou,

\textit{ETDataset: Electricity Transformer Dataset},

GitHub repository.

Available: \url{https://github.com/zhouhaoyi/ETDataset}


\bibitem{Yi2023}

K. Yi, Q. Zhang, W. Fan, S. Wang, P. Wang, H. He, N. An, D. Lian, L. Cao, and Z. Niu,

``Frequency-domain MLPs are More Effective Learners in Time Series Forecasting,''

in \textit{Proceedings of the Thirty-seventh Conference on Neural Information Processing Systems (NeurIPS)}, 2023.

Available: \url{https://openreview.net/forum?id=iif9mGCTfy}

\end{thebibliography}

\end{document}